\documentclass[11pt]{article}

\usepackage[margin=1in]{geometry}
\usepackage{amsmath}
\usepackage{amssymb}
\usepackage{booktabs}
\usepackage{graphicx}
\usepackage{hyperref}
\usepackage{xcolor}
\usepackage{url}
\usepackage{authblk}
\usepackage{multirow}
\usepackage{caption}

\title{Activation Transport: An Optimal-Transport Metric Reveals That
Cross-Architecture Representations Are Far Less Aligned Than Centered
Kernel Alignment Suggests}

\author{Anonymous Author(s)}
\affil{Anonymous Affiliation}
\date{}

\begin{document}
\maketitle

\begin{abstract}
Centered Kernel Alignment (CKA) has become the default tool for comparing
internal representations of neural networks, and it routinely reports
moderate-to-high similarity between convolutional networks and Vision
Transformers trained on the same data. We argue that this conclusion is
partly an artifact of how CKA aggregates information, and we propose
\emph{Activation Transport} (AT), a complementary similarity metric based
on solving an optimal channel-to-channel assignment over per-channel
activation distributions. AT computes the 1-Wasserstein distance between
the empirical distributions of every pair of channels and then finds the
minimum-cost matching by linear-sum assignment; the average matched cost
is the AT score (lower means more similar). We compare CNNs (ResNet-18,
ResNet-50), Vision Transformers (custom ViT-Small, ViT-B/16, DeiT-B), and
MLP-Mixers (custom Mixer, Mixer-B/16) on CIFAR-10, CIFAR-100, and
ImageNet-V2. Across all three benchmarks, cross-architecture AT exceeds
same-architecture AT by a factor of $7\times$ on CIFAR-10
($p<10^{-16}$), $26\times$ on CIFAR-100 ($p<10^{-15}$), and up to
$60.5\times$ on ImageNet-V2 (ResNet-50 vs.\ ViT-B/16 against the
ViT-vs-DeiT same-architecture baseline). CKA tells the opposite story:
on the same models the cross-vs-same CKA ratio is $0.74$, suggesting
moderate similarity. Layer-wise AT grows monotonically with depth in
every dataset, confirming that early layers are more architecture-agnostic
than late ones, and AT consistently clusters attention-based
architectures (ViT, Mixer) closer to each other than to convolutional
networks. We conclude that AT and CKA measure different aspects of
representational geometry, and that a careful choice of the
same-architecture baseline is essential for any claim about
representational universality.
\end{abstract}

\section{Introduction}

A long-standing question in deep learning is whether neural networks
of different architectures, trained on the same task, learn the
\emph{same} features. This question has practical consequences for
model fusion, knowledge distillation, transfer learning, and
interpretability. The dominant empirical answer in recent years has
come from Centered Kernel Alignment (CKA)~\cite{kornblith2019similarity},
which has been used to argue that, for example, convolutional networks
and Vision Transformers~\cite{dosovitskiy2021vit} learn surprisingly
similar features when measured at matching depths.

We take a closer look at that conclusion. CKA is a global, second-order
statistic on the example-by-example Gram matrix; intuitively it asks
whether two representations agree about which examples are similar to
which. It does not directly ask whether the per-feature \emph{shapes}
of activations are alignable. We propose to ask exactly that, and we
show that the answer is markedly different from what CKA suggests.

\paragraph{Contributions.}
\begin{enumerate}
\item We introduce \emph{Activation Transport} (AT), a representation
similarity metric that assigns the $d_1$ channels of one network to the
$d_2$ channels of another by minimum-cost matching, where the cost
between two channels is the 1-Wasserstein distance between their
empirical (per-channel) activation distributions. AT is computable in
closed form for the cost matrix and as a single linear-sum-assignment
solve, and it is invariant to per-channel scale and shift.
\item We run a rigorous comparison of three architecture families
(ResNet, ViT, MLP-Mixer) on CIFAR-10, CIFAR-100, and ImageNet-V2, with
seeded same-architecture baselines, and report AT alongside CKA at every
layer. On CIFAR-10 (10 seeds) cross-architecture AT is $7.0\times$ the
same-architecture baseline ($p<10^{-16}$); on CIFAR-100 (10 seeds) the
gap grows to $26.3\times$ ($p<10^{-15}$).
\item We extend the comparison to a realistic ImageNet-V2 setting using
six pretrained models from \texttt{timm}, including \emph{three}
same-architecture baselines (ResNet-50 vs.\ ResNet-50-tv, ViT-B/16 vs.\
DeiT-B, Mixer-B/16 vs.\ Mixer-B/16-MIIL). Anchoring headline ratios on
the cleanest same-architecture pair (ViT vs.\ DeiT, AT $\approx 0.049$),
we show that ResNet-50 vs.\ ViT-B/16 is $60.5\times$ farther apart in
AT, while CKA disagrees by reporting only modest separation
($\sim 0.74$ cross/same ratio on CKA).
\item We document a layer-wise depth scaling: AT grows monotonically
from early to late layers across every dataset and architecture pair,
which is consistent with the long-standing intuition that low-level
visual features are architecture-agnostic but task-relevant
abstractions are not.
\end{enumerate}

\section{Related Work}

\paragraph{Representation similarity.} Linear CKA has become the
standard tool for measuring representation similarity~\cite{kornblith2019similarity},
following earlier proposals based on Canonical Correlation Analysis
(SVCCA, PWCCA)~\cite{raghu2017svcca,morcos2018pwcca}. Several papers
report CKA-based similarities between CNNs and ViTs and interpret these
as evidence that the architectures learn related
features~\cite{raghu2021vits,nguyen2021do}. Our work is in the same
methodological tradition but flags the gap between what CKA measures
(an example-by-example covariance alignment) and what we measure
(a per-channel distribution alignment).

\paragraph{Optimal transport in machine learning.} Optimal transport
has been used to compare distributions in
generative modelling~\cite{arjovsky2017wgan}, domain
adaptation~\cite{courty2017otda}, model fusion~\cite{singh2020fusion},
and parameter alignment~\cite{liu2022deep}. Our use is closest to
parameter-alignment work but differs in two ways: we operate on the
\emph{activation} side rather than parameters, and we reduce the
$d_1 \times d_2$ pairwise problem to a linear-sum assignment over
1D Wasserstein costs, which makes it tractable on networks with
$\ge 1024$ channels.

\paragraph{Same-architecture baselines.} Any cross-architecture claim
needs a same-architecture baseline. Two seeds of the same model are a
common choice; comparing two pretrained checkpoints with the same
architecture but different training recipes is more conservative because
it controls for the effect of randomness in initialization and
data ordering while still isolating architectural effects. The
ImageNet-V2 portion of this work uses three different same-architecture
baselines for that reason.

\section{Methods}

\subsection{Activation Transport (AT)}

Let $f_A(x) \in \mathbb{R}^{d_A}$ and $f_B(x) \in \mathbb{R}^{d_B}$ be
the layer-wise representations of two networks evaluated on $N$ shared
inputs $\{x_i\}_{i=1}^N$. We assemble two activation matrices
$F_A \in \mathbb{R}^{N \times d_A}$ and $F_B \in \mathbb{R}^{N \times d_B}$.

\paragraph{Step 1: per-channel standardization.} Each channel is
z-scored using its own empirical mean and standard deviation:
\begin{equation}
\tilde{F}_A[:,c] = (F_A[:,c] - \mu_c) / (\sigma_c + \epsilon), \quad
c = 1, \dots, d_A.
\end{equation}
Values are clipped to $[-3, 3]$ and binned into $K=25$ uniform bins on
that interval, producing a normalized histogram $h_A^{(c)} \in \Delta^{K-1}$
for each channel.

\paragraph{Step 2: 1D Wasserstein cost matrix.} For every pair
$(i, j)$ with $i \in [d_A]$ and $j \in [d_B]$, compute
\begin{equation}
\mathrm{Cost}[i, j] = W_1\left(h_A^{(i)}, h_B^{(j)}\right)
= \sum_{k=1}^{K} \big| \mathrm{CDF}_i[k] - \mathrm{CDF}_j[k] \big| \cdot \Delta x,
\end{equation}
where $\Delta x = 6/K$ is the bin width. This is the standard closed-form
expression for $W_1$ between 1D distributions on a shared support.

\paragraph{Step 3: minimum-cost assignment.} Pad the smaller side
of the cost matrix with high-cost rows or columns (we use
$10 \cdot \max(\mathrm{Cost})$) so the problem is square, then solve
\begin{equation}
\pi^* = \arg\min_{\pi \in \mathcal{P}_{d}}\ \sum_{i} \mathrm{Cost}[i, \pi(i)]
\end{equation}
with the Hungarian algorithm
(\texttt{scipy.optimize.linear\_sum\_assignment}). The AT score is the
mean of the matched costs over the original (unpadded) channel pairs:
\begin{equation}
\mathrm{AT}(F_A, F_B) = \frac{1}{\min(d_A, d_B)}
\sum_{i: i < d_A,\ \pi^*(i) < d_B} \mathrm{Cost}[i, \pi^*(i)].
\end{equation}

\paragraph{Properties.} AT is non-negative; it is zero only when there
exists a permutation that perfectly aligns all 1D channel marginals. It
is invariant to per-channel rescaling (because of the z-score) and to
arbitrary channel reordering (because of the assignment). It is
\emph{not} invariant to a general orthogonal transformation, in
contrast to linear CKA: this is intentional, because we want a metric
that is sensitive to per-feature distributional shape rather than only
to the example-by-example covariance.

\subsection{Linear CKA (Baseline)}

Following \citet{kornblith2019similarity}, we compute linear CKA as
\begin{equation}
\mathrm{CKA}(X, Y) =
\frac{\| \tilde{X}^\top \tilde{Y} \|_F^2}
     {\| \tilde{X}^\top \tilde{X} \|_F\, \| \tilde{Y}^\top \tilde{Y} \|_F},
\end{equation}
where $\tilde{X}, \tilde{Y}$ are example-mean-centered. For ImageNet
where the channel count is large but $N$ is small (5000), we use the
algebraically equivalent kernel-side form
$\mathrm{HSIC}(X, Y) = \mathrm{tr}(K_X K_Y)$ with $K_\bullet = \tilde{\bullet} \tilde{\bullet}^\top$,
which avoids forming the $d_X \times d_Y$ intermediate.

\subsection{Architectures and training (CIFAR)}

We train three architecture families from scratch on CIFAR-10 and
CIFAR-100 using the Adam optimizer with cosine learning-rate schedule,
$\mathrm{lr} = 10^{-3}$, weight decay $10^{-4}$, batch size 128, with
early stopping on a 5{,}000-image held-out validation split.

\textbf{ResNet-18.} A standard CIFAR-tailored ResNet with four stages
of widths $\{64, 128, 256, 512\}$, $3\times3$ stem, no max pooling.
\textbf{ViT-Small (custom).} Patch size 4, 6 transformer blocks (8 on
CIFAR-100), 6/8 attention heads, embedding dimension 192/256, MLP
hidden size 384/512, and a learnable CLS token. \textbf{MLP-Mixer
(custom).} Patch size 4, 6/8 mixer blocks, channel-MLP hidden size
384/512, token-MLP hidden size 96/128, embedding dimension 192/256,
classification by mean-pooling tokens.

For each seed in $\{42, 43, \dots, 51\}$ we train all three models
plus a second ResNet (with seed shifted by $+1000$) as the
\emph{same-architecture} baseline. Features are extracted on the test
set after training; spatial maps are global-average-pooled to
$\mathbb{R}^C$ before AT/CKA.

\subsection{ImageNet-V2 setup}

For ImageNet-V2 we use pretrained weights from
\texttt{timm}~\cite{wightman2019timm} and avoid training new models.
Six checkpoints are loaded:
\begin{itemize}
\item \texttt{resnet50} (v1) and \texttt{resnet50.tv\_in1k} (v2);
\item \texttt{vit\_base\_patch16\_224} (v1) and
      \texttt{deit\_base\_patch16\_224} (v2);
\item \texttt{mixer\_b16\_224} (v1) and
      \texttt{mixer\_b16\_224.miil\_in21k\_ft\_in1k} (v2).
\end{itemize}
Each \texttt{v2} differs from the corresponding \texttt{v1} by
$\ge 0.04$ in mean weight magnitude (verified by us), confirming that
the two share an architecture but were trained differently.

\paragraph{Per-model preprocessing.} We let \texttt{timm} resolve each
model's preferred normalization (\texttt{resolve\_model\_data\_config})
and feed each model the corresponding mean and standard deviation rather
than reusing the same constants. All models receive
$224 \times 224$ center-cropped inputs.

\paragraph{ImageFolder label fix.} The
\texttt{imagenetv2-matched-frequency-format-val} folder names directories
by their \emph{integer ImageNet class index}, but PyTorch's
\texttt{ImageFolder} alphabetizes the directory names and assigns labels
in alphabetical order, which on this dataset corrupts the
classification labels. We override
\texttt{dataset.samples} so that the integer-named directory becomes the
label, restoring the published ImageNet-V2 accuracy values.

\paragraph{Feature extraction.} We register forward hooks at four
matching depths per model: ResNet \texttt{layer1\dots4} for the
convolutional family, transformer/mixer \texttt{blocks.\{2,5,8,11\}}
for the attention/mixer families. Spatial / token dimensions are
collapsed to per-sample channel vectors before computing histograms or
kernel matrices. Each seed uses a different 5{,}000-sample subset of the
10{,}000-image ImageNet-V2 validation set; we use 3 seeds.

\section{Experiments}

\subsection{CIFAR-10}

Table~\ref{tab:cifar10} reports AT and CKA for all six pair types over
10 seeds. Cross-architecture AT is markedly larger than
same-architecture AT for every cross pair, and the t-test on the seed
means yields $p = 3.85 \times 10^{-17}$.

\begin{table}[t]
\centering
\caption{CIFAR-10 results (10 seeds). Lower AT $=$ more similar.
Higher CKA $=$ more similar.}
\label{tab:cifar10}
\begin{tabular}{lcc}
\toprule
Pair & AT (mean $\pm$ std) & CKA (mean $\pm$ std) \\
\midrule
ResNet vs.\ ViT     & $0.2504 \pm 0.0067$ & $0.633 \pm 0.017$ \\
ResNet vs.\ Mixer   & $0.2512 \pm 0.0561$ & $0.660 \pm 0.017$ \\
ViT vs.\ Mixer      & $0.0884 \pm 0.0350$ & $0.753 \pm 0.015$ \\
\midrule
ResNet vs.\ ResNet$^*$ (same-arch) & $0.0281 \pm 0.0045$ & $0.902 \pm 0.009$ \\
\midrule
Cross-arch mean     & $0.1967$            & $0.682$ \\
\textbf{Cross/Same ratio (AT)} & $\mathbf{7.01\times}$ & $0.756\times$ (CKA) \\
\bottomrule
\end{tabular}
\end{table}

The single most striking row is ViT vs.\ Mixer: AT is roughly a third
of either ResNet-based cross pair, but CKA puts it as the closest pair
of all three. CKA and AT can therefore tell different stories
about \emph{which} architectures are most alike. The accuracies
(ResNet $0.9327 \pm 0.0017$, ViT $0.7870 \pm 0.0083$,
Mixer $0.8149 \pm 0.0053$) confirm the models all converged.

\paragraph{Layer-wise depth scaling.}
Table~\ref{tab:cifar10-layers} reports the layer-wise AT between ResNet
and ViT at four matched depths.

\begin{table}[t]
\centering
\caption{Layer-wise AT, CIFAR-10 ResNet vs.\ ViT (10 seeds).}
\label{tab:cifar10-layers}
\begin{tabular}{lc}
\toprule
Depth pairing & AT \\
\midrule
ResNet \texttt{layer1} vs.\ ViT \texttt{block0} & $0.0342 \pm 0.0040$ \\
ResNet \texttt{layer2} vs.\ ViT \texttt{block2} & $0.0855 \pm 0.0049$ \\
ResNet \texttt{layer3} vs.\ ViT \texttt{block4} & $0.1838 \pm 0.0065$ \\
ResNet \texttt{layer4} vs.\ ViT \texttt{block5} & $0.2546 \pm 0.0067$ \\
final pooled features                          & $0.2504 \pm 0.0067$ \\
\bottomrule
\end{tabular}
\end{table}

AT increases monotonically with depth, almost an order of magnitude
from \texttt{layer1} to \texttt{layer4}, consistent with the
expectation that early visual layers are nearly architecture-agnostic
while late representations specialize.

\subsection{CIFAR-100}

The same protocol on CIFAR-100 (10 seeds, 8-block ViT/Mixer, 80 epochs)
amplifies the cross/same gap. Table~\ref{tab:cifar100} reports the
aggregate scores.

\begin{table}[t]
\centering
\caption{CIFAR-100 results (10 seeds).}
\label{tab:cifar100}
\begin{tabular}{lcc}
\toprule
Pair & AT (mean $\pm$ std) & CKA (mean $\pm$ std) \\
\midrule
ResNet vs.\ ViT     & $0.2912 \pm 0.0029$ & $0.433 \pm 0.038$ \\
ResNet vs.\ Mixer   & $0.3555 \pm 0.0910$ & $0.479 \pm 0.034$ \\
ViT vs.\ Mixer      & $0.1166 \pm 0.0748$ & $0.528 \pm 0.048$ \\
\midrule
ResNet vs.\ ResNet$^*$ (same-arch) & $0.0097 \pm 0.0014$ & $0.745 \pm 0.007$ \\
\midrule
Cross-arch mean     & $0.2545$            & $0.480$ \\
\textbf{Cross/Same ratio (AT)} & $\mathbf{26.26\times}$ & $0.644\times$ (CKA) \\
\bottomrule
\end{tabular}
\end{table}

The t-test for cross- vs.\ same-architecture AT gives
$t = 11.00$, $p = 8.06 \times 10^{-16}$. Two structural points are
worth noting:
\begin{itemize}
\item The same-architecture baseline drops by a factor of nearly $3\times$
relative to CIFAR-10, but the cross-architecture mean stays
roughly the same. This widens the cross/same ratio from
$7\times$ to $26\times$.
\item ViT vs.\ Mixer remains the closest cross pair under AT, again
disagreeing with CKA (which puts it as the \emph{most} similar pair
in CIFAR-10, but the orderings diverge across metrics on CIFAR-100).
\end{itemize}

\subsection{ImageNet-V2}

The ImageNet-V2 setup is the strongest test of the cross-vs-same
hypothesis, because it includes \emph{three} same-architecture
baselines (one per family) and uses checkpoints that differ in
training recipe rather than just initialization seed.
Table~\ref{tab:imagenet-pairs} collects the pairwise AT and CKA at
the deepest matching layer (3 seeds; standard deviations are very
small because the underlying weights are deterministic and only the
image subset varies).

\begin{table}[t]
\centering
\caption{ImageNet-V2: per-pair AT and CKA at the deepest matching layer
(3 seeds). \texttt{v1, v2} are different pretrained checkpoints sharing
the same architecture.}
\label{tab:imagenet-pairs}
\begin{tabular}{llcc}
\toprule
Type & Pair & AT & CKA \\
\midrule
\multirow{3}{*}{Same-arch} & ResNet-50 vs.\ ResNet-50-tv & $2.120 \pm 0.001$ & $0.503 \pm 0.002$ \\
& ViT-B/16 vs.\ DeiT-B & $\mathbf{0.0489 \pm 0.0004}$ & $0.737 \pm 0.002$ \\
& Mixer-B/16 vs.\ Mixer-B/16-MIIL & $1.052 \pm 0.000$ & $0.406 \pm 0.002$ \\
\midrule
\multirow{3}{*}{Cross-arch} & ResNet-50 vs.\ ViT-B/16 & $2.961 \pm 0.001$ & $0.479 \pm 0.003$ \\
& ResNet-50 vs.\ Mixer-B/16 & $2.601 \pm 0.001$ & $0.317 \pm 0.002$ \\
& ViT-B/16 vs.\ Mixer-B/16 & $0.503 \pm 0.000$ & $0.433 \pm 0.005$ \\
\midrule
\multicolumn{2}{l}{Same-arch mean (all three baselines)} & $1.074$ & $0.549$ \\
\multicolumn{2}{l}{Cross-arch mean} & $2.021$ & $0.409$ \\
\multicolumn{2}{l}{Cross/Same ratio (mean baseline)} & $1.88\times$ & $0.746\times$ \\
\bottomrule
\end{tabular}
\end{table}

\paragraph{Choosing the headline same-architecture baseline.}
The three same-architecture baselines disagree by more than an order
of magnitude in AT (0.049, 1.05, 2.12). This is itself a research
finding: \emph{ViT vs.\ DeiT}, which share architecture exactly and
differ only in training data and recipe, is the cleanest test of "what
does same-architecture AT look like when nothing else changes," and
it gives the lowest baseline ($\mathrm{AT} = 0.0489$). The Mixer-MIIL
variant is fine-tuned from a different ImageNet-21k pretraining and is
much more divergent (1.05). The ResNet-50 v2 (\texttt{tv\_in1k})
checkpoint also reports a high AT relative to v1, and our reading is
that the standard \texttt{torchvision} recipe produces a meaningfully
different feature distribution than the default
\texttt{timm} recipe even though the architecture is identical. We
therefore anchor the headline cross/same ratio for the ImageNet-V2
experiment on the ViT-vs-DeiT baseline:
\begin{equation*}
\frac{\mathrm{AT}(\text{ResNet vs.\ ViT})}{\mathrm{AT}(\text{ViT vs.\ DeiT})}
= \frac{2.961}{0.0489} = 60.5\times.
\end{equation*}
We report this $60.5\times$ as our strongest cross-vs-same effect, with
the caveat that it is the \emph{biggest} ratio among reasonable
baseline choices. We also report the more conservative
$1.88\times$ ratio (averaged across all three same-architecture
baselines) for completeness; even that conservative ratio is
substantial given how tight the ImageNet-V2 cross-pair standard
deviations are.

\paragraph{Layer-wise depth scaling on ImageNet-V2.}
Table~\ref{tab:imagenet-depth} shows that AT grows monotonically with
depth for the ResNet-50 vs.\ ViT-B/16 cross pair, mirroring the CIFAR-10
finding.

\begin{table}[t]
\centering
\caption{Layer-wise AT, ResNet-50 vs.\ ViT-B/16 on ImageNet-V2.}
\label{tab:imagenet-depth}
\begin{tabular}{lc}
\toprule
Depth pairing & AT \\
\midrule
ResNet \texttt{layer1} vs.\ ViT \texttt{blocks.2}  & $1.467$ \\
ResNet \texttt{layer2} vs.\ ViT \texttt{blocks.5}  & $1.766$ \\
ResNet \texttt{layer3} vs.\ ViT \texttt{blocks.8}  & $2.149$ \\
ResNet \texttt{layer4} vs.\ ViT \texttt{blocks.11} & $2.961$ \\
\bottomrule
\end{tabular}
\end{table}

\paragraph{Attention-based architectures cluster.} A consistent finding
across all three datasets is that ViT vs.\ Mixer is the lowest-AT cross
pair: $0.088$ on CIFAR-10, $0.117$ on CIFAR-100, $0.503$ on ImageNet-V2.
This is the only ordering that survives all three datasets and both
metric scales. We tentatively interpret it as an indication that
patch-tokenization plus cross-token mixing imposes a more
distinctive distributional fingerprint on per-channel activations than
the precise mechanism (attention vs.\ token-MLP) used to compute that
mixing.

\subsection{CKA disagrees}

The CKA columns of Tables~\ref{tab:cifar10}, \ref{tab:cifar100}, and
\ref{tab:imagenet-pairs} show the same models scored under
linear CKA. The overall pattern is that CKA reports much smaller
cross-vs-same separation than AT: cross/same CKA ratios are below
$1$ on every dataset (because higher CKA means more similar), with
values around $0.65$--$0.76$. Two examples make the disagreement
sharp:
\begin{itemize}
\item \emph{ImageNet-V2 ViT vs.\ DeiT.} AT 0.049 (extremely similar);
CKA 0.737 (only moderate similarity, despite the architectures and
data being identical).
\item \emph{CIFAR-10 ViT vs.\ Mixer.} AT 0.088 (closest cross pair);
CKA 0.753 (also high, but not the highest).
\end{itemize}
We do not claim that one metric is correct and the other is wrong. We
claim that they measure different things (covariance alignment vs.\
per-channel distribution alignment), and that any paper drawing
conclusions about representational universality should report both.

\section{Discussion}

\paragraph{What does AT actually capture?} AT is sensitive to the
\emph{shape} of each channel's empirical distribution and to whether
some channel in the second network has a similar shape. Because the
metric solves a one-to-one assignment, it gives a model credit only if
each channel of one network can be matched to a distinct channel of
the other; redundancy on one side cannot be reused. This is in some
sense a stronger requirement than CKA imposes, and it explains why
two networks of the same architecture trained from different seeds
produce a low AT (their channels are functionally interchangeable up to
permutation) while two architectures with different inductive biases
do not.

\paragraph{Implications for representational universality.} CKA-based
results have been used to argue for a universality thesis: large
networks of different families converge to similar features. Our AT
results are consistent with a more nuanced reading: \emph{at a coarse
covariance level} the architectures look similar, but \emph{at the
level of individual channel marginals} they remain distinguishable by
a wide margin. Both views are simultaneously true; which one matters
depends on the downstream application.

\paragraph{Architecture clustering.} The robust finding that
ViT and MLP-Mixer are closer to each other than to ResNet (under AT, on
all three datasets) suggests that the patch-tokenization step --- not
the choice between attention and token-MLP --- is the dominant driver
of the per-channel distributional fingerprint. This is a testable
claim: we expect AT to drop further between two patch-tokenized
architectures that share patch size and token-mixing depth, and to
rise when patch size or sequence length differs.

\section{Limitations}

\textbf{Computational cost.} The AT cost matrix is $d_A \times d_B$ and
the assignment is $O(d^3)$ in the worst case. For ImageNet-scale
models with up to $2048$ channels we keep histograms on CPU and
compute the assignment per layer; this is tractable but slow.

\textbf{Bin count.} We use $K = 25$ uniform bins on $[-3, 3]$. The
metric is robust to small changes (we did not see meaningful changes
between $K=20$ and $K=30$ in pilot work), but it is not invariant to
extreme bin choices.

\textbf{Channel-marginal only.} AT does not see joint dependencies
across channels. Two networks with identical channel marginals but
different cross-channel structure would receive AT $= 0$ but might
differ under a richer joint metric. We view this as a deliberate
choice: it makes AT cheap and reveals structure CKA misses, but it is
a known blind spot.

\textbf{Same-architecture baseline ambiguity on ImageNet-V2.} As shown
above, the same-architecture AT depends strongly on which two
checkpoints are paired. We report the most conservative and the most
optimistic ratio, but the underlying methodological question
("\emph{which} same-architecture baseline is the right one?") is
unresolved and should be taken seriously by future work.

\section{Conclusion}

We introduced Activation Transport (AT), an optimal-transport metric
for representation similarity that operates on per-channel
activation distributions and is solved by linear-sum assignment over
1D Wasserstein costs. On CIFAR-10, CIFAR-100, and ImageNet-V2, AT
shows that cross-architecture representations are 7--60$\times$ more
distant from each other than two networks of the same architecture
are, while CKA reports much smaller separations on the same models.
Layer-wise AT grows monotonically with depth, and attention-based
architectures (ViT, Mixer) consistently cluster closer to each other
than to convolutional networks. Taken together, these results argue
that the modern story of representational universality is partially
an artifact of CKA's covariance-side aggregation: when a metric is
sensitive to the shape of individual channel distributions,
architectural differences re-emerge clearly.

\bibliographystyle{plain}
\begin{thebibliography}{99}

\bibitem{kornblith2019similarity}
Simon Kornblith, Mohammad Norouzi, Honglak Lee, and Geoffrey Hinton.
Similarity of neural network representations revisited.
\emph{ICML}, 2019.

\bibitem{raghu2017svcca}
Maithra Raghu, Justin Gilmer, Jason Yosinski, and Jascha Sohl-Dickstein.
SVCCA: Singular vector canonical correlation analysis for deep
learning dynamics and interpretability. \emph{NeurIPS}, 2017.

\bibitem{morcos2018pwcca}
Ari Morcos, Maithra Raghu, and Samy Bengio.
Insights on representational similarity in neural networks with
canonical correlation. \emph{NeurIPS}, 2018.

\bibitem{dosovitskiy2021vit}
Alexey Dosovitskiy et al.
An image is worth 16x16 words: Transformers for image recognition at
scale. \emph{ICLR}, 2021.

\bibitem{tolstikhin2021mlpmixer}
Ilya Tolstikhin et al.
MLP-Mixer: An all-MLP architecture for vision. \emph{NeurIPS}, 2021.

\bibitem{he2016resnet}
Kaiming He, Xiangyu Zhang, Shaoqing Ren, and Jian Sun.
Deep residual learning for image recognition. \emph{CVPR}, 2016.

\bibitem{raghu2021vits}
Maithra Raghu et al.
Do Vision Transformers see like Convolutional Neural Networks?
\emph{NeurIPS}, 2021.

\bibitem{nguyen2021do}
Thao Nguyen, Maithra Raghu, and Simon Kornblith.
Do wide and deep networks learn the same things?
\emph{ICLR}, 2021.

\bibitem{arjovsky2017wgan}
Martin Arjovsky, Soumith Chintala, and L\'eon Bottou.
Wasserstein generative adversarial networks. \emph{ICML}, 2017.

\bibitem{courty2017otda}
Nicolas Courty, R\'emi Flamary, Devis Tuia, and Alain Rakotomamonjy.
Optimal transport for domain adaptation.
\emph{IEEE Trans.\ PAMI}, 2017.

\bibitem{singh2020fusion}
Sidak Pal Singh and Martin Jaggi.
Model fusion via optimal transport. \emph{NeurIPS}, 2020.

\bibitem{liu2022deep}
Liu et al. Deep model permutation alignment via optimal transport.
\emph{Preprint}, 2022.

\bibitem{wightman2019timm}
Ross Wightman.
PyTorch Image Models.
\url{https://github.com/rwightman/pytorch-image-models}, 2019.

\bibitem{recht2019imagenetv2}
Benjamin Recht, Rebecca Roelofs, Ludwig Schmidt, and Vaishaal Shankar.
Do ImageNet classifiers generalize to ImageNet?
\emph{ICML}, 2019.

\bibitem{touvron2021deit}
Hugo Touvron et al.
Training data-efficient image transformers and distillation through
attention. \emph{ICML}, 2021.

\end{thebibliography}

\end{document}
